% Generated draft: requires mathematical and source review.
\documentclass[11pt]{article}
\usepackage[margin=25mm]{geometry}
\usepackage{fontspec}
\setmainfont{Latin Modern Roman}
\usepackage{amsmath,amssymb,mathtools,unicode-math}
\setmathfont{Latin Modern Math}
\usepackage{enumitem,needspace,xurl,hyperref}
\setlength{\parindent}{0pt}
\setlength{\parskip}{5pt}
\setlength{\emergencystretch}{3em}
\newcommand{\N}{\mathbb{N}}
\newcommand{\Z}{\mathbb{Z}}
\newcommand{\Q}{\mathbb{Q}}
\newcommand{\R}{\mathbb{R}}
\newcommand{\C}{\mathbb{C}}
\newcommand{\F}{\mathbb{F}}
\newcommand{\A}{\mathbb{A}}
\newcommand{\PP}{\mathbb{P}}
\newcommand{\E}{\mathbb{E}}
\newcommand{\eps}{\varepsilon}
\newcommand{\dd}{\,\mathrm{d}}
\begin{document}
\section*{UnsolvedMath ID: 800652}
\textbf{Source problem code:} catalogue entry
\subsection{Definitions and mathematical statement}
The supplied question asks about three splitting properties. Its statement does not specify the ambient class of spaces or define the cover classes. The following is an explicit conventional interpretation, conditional on confirming those conventions in the source.

Let \(X\) be a topological space. For this interpretation, a cover is a countably infinite family \(\mathcal U\) of open subsets of \(X\) such that \(\bigcup\mathcal U=X\) and \(X\notin\mathcal U\). Let \(\Omega(X)\) consist of those covers \(\mathcal U\) for which
\[
\forall F\subseteq X\;\bigl(F\text{ finite}\Longrightarrow \exists U\in\mathcal U\;(F\subseteq U)\bigr).
\]
A cover \(\mathcal U\) is large if every point belongs to infinitely many members of \(\mathcal U\). Let \(\mathrm T(X)\) consist of the large covers satisfying, for every \(x,y\in X\),
\[
\bigl|\{U\in\mathcal U:x\in U,\ y\notin U\}\bigr|<\infty
\quad\text{or}\quad
\bigl|\{U\in\mathcal U:y\in U,\ x\notin U\}\bigr|<\infty.
\]
For cover classes \(\mathcal A(X)\) and \(\mathcal B(X)\), say that \(X\) satisfies \(\mathsf{Split}(\mathcal A,\mathcal B)\) if every \(\mathcal U\in\mathcal A(X)\) admits a partition
\[
\mathcal U=\mathcal U_0\mathbin{\dot\cup}\mathcal U_1,
\qquad \mathcal U_0,\mathcal U_1\in\mathcal B(X).
\]

For each of the three properties
\[
P\in\{\mathsf{Split}(\Omega,\Omega),\ \mathsf{Split}(\Omega,\mathrm T),\ \mathsf{Split}(\mathrm T,\mathrm T)\},
\]
the question is whether the following preservation assertion holds:
\[
\forall X\in\mathcal C\;\forall n\in\mathbb Z_{\geq 1},
\qquad X\models P\ \Longrightarrow\ X^n\models P,
\]
where \(X^n\) carries the product topology and \(\mathcal C\) denotes the intended ambient class of spaces. The supplied record leaves \(\mathcal C\) unspecified and does not establish the countability and openness conventions used above. Consequently, this is a conditional formalization rather than a fully source-verified formulation. Each of the three preservation assertions is a separate target; the wording does not assert that any one of them holds.
\subsection{Short English statement}
For each of the three stated splitting properties, ask whether every positive finite Cartesian power of a space with that property also has it. The intended class of spaces and the precise cover conventions require source verification. The registry labels the problem open; that label is supplied metadata, not an independently verified status assessment.
\subsection{Sources}
\begin{itemize}
\item Local UnsolvedMath registry, permanent ID \(800652\), exact source problem code \texttt{catalogue entry}. The supplied record attributes the problem to Tsaban, gives the proposed year as 2003, and labels its status \texttt{open}.
\item Supplied source list: Tsaban, \emph{Selection Principles in Mathematics: A Milestone of Open Problems} (2003), Problem 25. Supplied URL: \url{https://arxiv.org/abs/math/0312182}. The record reports access on 2026-10-16 and extraction by \texttt{arxiv\_latex}; these are supplied provenance details, not a claim of independent access or verification.
\item Supplied attribution/context: \texttt{[splittability, Problem 8.8]}. The record supplies no bibliographic expansion of the reference label \texttt{splittability}.
\end{itemize}
\end{document}
